\chapter{Modularidad, dominio y patrones}

\begin{objetivos}
El lector podrá diseñar límites modulares, reconocer cohesión y acoplamiento, emplear conceptos de diseño orientado al dominio y construir una arquitectura hexagonal que mantenga el núcleo independiente de interfaces y persistencia.
\end{objetivos}

\section{Modularidad antes que distribución}
Un módulo agrupa responsabilidades que cambian por razones relacionadas y expone un contrato deliberado. La modularidad es una propiedad lógica; no exige múltiples procesos. Un monolito puede estar bien modularizado y un conjunto de servicios puede formar un ``monolito distribuido'' si todos comparten base de datos, despliegue y conocimiento interno.

La cohesión es la relación interna entre los elementos de un módulo. El acoplamiento es la dependencia entre módulos. El objetivo no es eliminar dependencias, sino hacerlas necesarias, explícitas y estables. Una dependencia hacia una abstracción de negocio suele ser menos volátil que una dependencia hacia una biblioteca concreta de infraestructura.

Para LibreReserva se proponen cuatro contextos iniciales:
\begin{itemize}
  \item \textbf{Recursos}: salas, equipos, capacidades y restricciones;
  \item \textbf{Reservas}: disponibilidad, reglas de solapamiento y ciclo de vida;
  \item \textbf{Identidad}: usuarios, roles y relación con el proveedor institucional;
  \item \textbf{Notificaciones}: preferencias, plantillas y entrega asíncrona.
\end{itemize}

\section{Lenguaje ubicuo y límites de dominio}
El diseño orientado al dominio propone que desarrolladores y expertos compartan vocabulario. ``Reserva confirmada'' debe poseer el mismo significado en conversaciones, código, eventos y pruebas. Cuando una palabra cambia de sentido entre áreas, puede indicar una frontera. Por ejemplo, un \emph{recurso} para inventario puede incluir costo de reposición, mientras para reservas sólo interesa capacidad y disponibilidad.

Un agregado protege invariantes dentro de una frontera transaccional. No debe crecer para contener todo lo relacionado. En LibreReserva, la regla ``un recurso no puede tener dos reservas confirmadas que se solapen'' requiere una estrategia de concurrencia en Reservas; no obliga a cargar información completa del usuario ni el historial de notificaciones.

\section{Arquitectura hexagonal}

\begin{figure}[H]
\centering
\begin{tikzpicture}[every node/.style={align=center}]
  \node[draw=AzulUD,fill=AzulClaro,regular polygon,regular polygon sides=6,minimum size=4.2cm] (core) {Dominio y\\casos de uso};
  \node[draw=Verde,fill=VerdeClaro,rounded corners,left=2.5cm of core] (http) {Adaptador\\HTTP/FastAPI};
  \node[draw=Verde,fill=VerdeClaro,rounded corners,right=2.5cm of core] (db) {Adaptador\\PostgreSQL};
  \node[draw=Naranja,fill=NaranjaClaro,rounded corners,above=1.4cm of core] (cli) {Pruebas / CLI};
  \node[draw=Naranja,fill=NaranjaClaro,rounded corners,below=1.4cm of core] (mq) {RabbitMQ / correo};
  \draw[-{Latex},thick] (http) -- node[above]{puerto de entrada} (core);
  \draw[-{Latex},thick] (core) -- node[above]{puerto de salida} (db);
  \draw[-{Latex},thick] (cli) -- (core);
  \draw[-{Latex},thick] (core) -- (mq);
\end{tikzpicture}
\caption{El dominio depende de puertos; los adaptadores dependen del dominio.}
\label{fig:hexagonal}
\end{figure}

Los puertos expresan capacidades requeridas o ofrecidas. Un caso de uso puede depender de \software{RepositorioReservas} sin conocer SQLAlchemy. El adaptador PostgreSQL implementa ese puerto. Esta dirección permite probar el núcleo con dobles simples y reemplazar infraestructura sin modificar reglas de negocio.

\section{Patrones y límites}
\subsection{MVC, Repository y DAO}
MVC separa entrada, coordinación y presentación, pero no define toda la arquitectura. Repository representa una colección de objetos del dominio; DAO encapsula acceso a datos con una orientación más cercana a tablas o consultas. Usarlos como sinónimos suele producir capas sin significado.

\subsection{Inversión de dependencias}
Las políticas de alto nivel no deberían depender de detalles de bajo nivel. La inversión se concreta cuando el dominio define el contrato y el adaptador lo implementa. Crear una interfaz para cada clase no mejora automáticamente el diseño: la abstracción debe representar variabilidad o un límite real.

\section{Esqueleto modular en Python}
\begin{instalacion}
Active el entorno virtual creado en el capítulo 1 e instale herramientas de desarrollo:
\begin{lstlisting}[language=bash]
python -m pip install --upgrade pip
python -m pip install "fastapi[standard-no-fastapi-cloud-cli]" pytest
\end{lstlisting}
FastAPI y pytest son proyectos de código abierto. Registre versiones con \comando{python -m pip freeze > requirements.lock}.
\end{instalacion}

\begin{lstlisting}[language=Python]
# src/reservas/dominio.py
from dataclasses import dataclass
from datetime import datetime

@dataclass(frozen=True)
class Intervalo:
    inicio: datetime
    fin: datetime

    def __post_init__(self):
        if self.inicio >= self.fin:
            raise ValueError("el intervalo debe ser positivo")

    def se_solapa(self, otro: "Intervalo") -> bool:
        return self.inicio < otro.fin and otro.inicio < self.fin
\end{lstlisting}

\begin{lstlisting}[language=Python]
# tests/test_intervalo.py
from datetime import datetime
import pytest
from src.reservas.dominio import Intervalo

def test_rechaza_intervalo_vacio():
    instante = datetime(2026, 8, 6, 10, 0)
    with pytest.raises(ValueError):
        Intervalo(instante, instante)

def test_detecta_solapamiento():
    a = Intervalo(datetime(2026, 8, 6, 9), datetime(2026, 8, 6, 11))
    b = Intervalo(datetime(2026, 8, 6, 10), datetime(2026, 8, 6, 12))
    assert a.se_solapa(b)
\end{lstlisting}

\begin{validacion}
Ejecute \comando{pytest -q}. Las pruebas del dominio deben funcionar sin iniciar FastAPI, PostgreSQL o RabbitMQ. Esa independencia es evidencia del límite, no sólo una preferencia de carpetas.
\end{validacion}

\section{Ejercicios}
\begin{ejercicio}[title={Ejercicio 1. Mapa de contextos}]
Defina propietarios, lenguaje, datos y contratos para Recursos, Reservas, Identidad y Notificaciones. Señale relaciones cliente--proveedor y términos que cambian de significado entre contextos.
\end{ejercicio}

\begin{ejercicio}[title={Ejercicio 2. Prueba de frontera}]
Cree una prueba que falle si el paquete \software{src/reservas/dominio.py} importa FastAPI, SQLAlchemy o un cliente de RabbitMQ. Puede analizar el árbol sintáctico con el módulo estándar \software{ast}. Explique qué tipos de acoplamiento esta prueba no puede detectar.
\end{ejercicio}

\begin{ejercicio}[title={Ejercicio 3. Concurrencia}]
Proponga tres estrategias para impedir reservas simultáneas: bloqueo pesimista, control optimista y restricción en base de datos. Compare consistencia, rendimiento, experiencia del usuario y complejidad de recuperación.
\end{ejercicio}
